\documentclass[12pt]{article}
\usepackage[utf8]{inputenc}
\usepackage[T1]{fontenc}
\usepackage{hyperref}
\title{Yinggarda Language Resources for K-12 Curriculum Development}
\author{Compiled by Claude Sonnet 5.5 (claude-sonnet-5-5) for Andr\'as Kornai}
\date{}
\begin{document}
\maketitle

\section{Introduction}
Yinggarda (also Yingkarta, Inggarda; ISO 639-3 yia; Glottocode ying1247) is a dormant Pama-Nyungan (Kartu) language of the Gascoyne coast, Western Australia. Speaker estimates conflict (about 2 to 5, or none). It is not known to be taught in schools.

\section{Primary Educational Resources}
No school materials were found.

\section{Linguistic and Reference Materials}
\begin{itemize}
\item Austin, \emph{A dictionary of Yinggarda, Western Australia} (La Trobe University, 1992, 69 pages); library record. \url{https://cockburnlibrary.spydus.com/cgi-bin/spydus.exe/FULL/OPAC/BIBENQ/114780683/12214210,32}
\item Peter Austin's page on Gascoyne region languages. It mentions a draft Yinggarda-English dictionary based on Alan Dench's 1978 and 1980 materials, with publication hoped for in 2025 or 2026 (not confirmed). \url{https://peterkaustin.com/?p=75}
\item Yinggarda language (Wikipedia). \url{https://en.wikipedia.org/wiki/Yinggarda_language}
\item AUSTLANG entry W19, Yingkarta. \url{https://collection.aiatsis.gov.au/austlang/language/W19}
\item Endangered Languages Project, Yinggarda. \url{https://www.endangeredlanguages.com/lang/3198}
\end{itemize}
Alan Dench's 1998 sketch grammar was mentioned in a search summary, but no URL was seen, so it is not listed as a resource.

\section{Religious and Community Resources}
\begin{itemize}
\item Gascoyne Development Commission, Aboriginal Culture page, on the region's language groups. \url{https://gdc.wa.gov.au/our-region/aboriginal-culture.aspx}
\end{itemize}
The summary also mentions the Yamaji/Irra Wangga Language Centre, which has worked on Yinggarda since 1993, and the Gwoonwardu Mia cultural centre in Carnarvon. No URLs for them were seen.

\section{Digital and Audio Resources}
No audio or digital teaching resources were found.

\section{Curriculum Development Context}
A heritage and revitalization curriculum would have to build on Austin's dictionary and Dench's materials. It should be developed with Yamaji community and language centre involvement. Further searching is needed.

\end{document}